\subsection{First-Class Citizens and Higher-Order Functions}

A function object can travel through the program like any other object: assigned to variables, stored in collections, passed as arguments, and returned from calls. A higher-order function receives or returns another function.

\subsubsection{Functions as In-Memory Objects: Passing, Returning, and Storing Subroutine References inside Variables}

\begin{verbatim}
def say_doh():
    return "D'oh!"

fn = say_doh
print(f"fn is say_doh: {fn is say_doh}, fn(): {fn()}")

def apply_twice(fn, arg):
    return fn(fn(arg))

def add_one(x):
    return x + 1

print(apply_twice(add_one, 10))
\end{verbatim}

Possible output:

\begin{verbatim}
fn is say_doh: True, fn(): D'oh!
12
\end{verbatim}

Writing \texttt{fn} passes the object; writing \texttt{fn()} invokes \texttt{\_\_call\_\_}.

\subsubsection{Callbacks and Callback Chains: Decoupling Structural Execution Graphs}

Dispatch tables map keys to function objects. Selection happens at runtime by indexing the table and calling the retrieved callable.

\begin{verbatim}
DISPATCH = {
    "double": lambda x: x * 2,
    "square": lambda x: x * x,
}

def run(op, value):
    return DISPATCH[op](value)

print(run("double", 21))
print(run("square", 7))
\end{verbatim}

Possible output:

\begin{verbatim}
42
49
\end{verbatim}

Lambdas in dispatch tables compile to ordinary \texttt{PyCodeObject} instances with the same cost profile as \texttt{def}.

\subsubsection{Higher-Order Functions: Functions That Receive or Return Other Functions}

\begin{verbatim}
def make_prefixer(prefix):
    def add_prefix(text):
        return f"{prefix}: {text}"
    return add_prefix

warn = make_prefixer("WARNING")
print(warn("System alert"))
\end{verbatim}

Possible output:

\begin{verbatim}
WARNING: System alert
\end{verbatim}

A higher-order function returns a configured callable whose closure cells hold state without globals or class instances.
